\documentclass[11pt]{article}
\usepackage{amsmath,amssymb,amsthm,hyperref}
\newtheorem{proposition}{Proposition}
\newcommand{\E}{\mathbb E}
\newcommand{\Pp}{\mathcal P}
\newcommand{\op}{\mathrm{op}}
\title{Centered nuclear duality and a sparse two-copy defect budget}
\author{Campaign 3395: Engineer continuation}
\date{November 4, 2026}
\begin{document}
\maketitle

\paragraph{Status and the missing result.}
The fixed-factor theorem, T5, and the signed dual charge below remain
unproved. The new results are the requested centering bridge, an exact
sparse-noise coupling with unchanged original one-hot marginals and one
bounded whole-output defect budget, and a high-frequency reduction with
a dimension-independent noise inverse. The budget has not been shown
to control the signed pairing. These are partial lemmas, not an
achieved fixed approximation factor or a publication qualification.
Reviewed V14 remains the strongest achieved approximation theorem.

Reuse independently reviewed Z1--Z20, V1--V16, T1--T4, J8--J12,
K1--K5, the original sampler and cut, and A4/A7 without unchanged
replay. Nuclear/operator duality is already independently attributed
to arXiv:0706.4138v1, Proposition 2.1 and equations (2.5)--(2.6).
It is used only as a finite-dimensional matrix identity; a short
SVD argument is also given below. No matrix-recovery or circuit
stability theorem, new external assumption, or novelty claim is used.
The capped relative-mass route remains retired for its reviewed reason.

\section{The elementary bridge and its exact dual}
Retain every definition and conditioning in Z1. In particular
$E'_j=\min(R_j,3)-M_j(Q_j)$, $0\le E'_j\le3$,
$m_j=2^{k_j}$, and $a_j$ is the conditional mean of the actual
heavy part, not the mean of the error. Write
\[
 I_j=1_{\{k_j\ge5k_0/4,\ a_j\ge2^{-k_0}\}},\quad
 P_j=\frac{\boldsymbol1\boldsymbol1^{\mathsf T}}{m_j},\quad
 X_j=(I-P_j)E'_j(I-P_j),\qquad k_0=N/2^{134}.
\]
All energies are zero outside $I_j=1$. Put
$\mathfrak N_{\rm cen}=\sum_j\E[I_j\|X_j\|_*^2/(m_j^3a_j)]$.
The expectation is exactly Z1's original law.

\begin{proposition}
For the unrestricted circuits in Z1,
\[
 \mathfrak N_{\rm cen}\le\mathfrak N
 \le2\mathfrak N_{\rm cen}+36s\,2^{-k_0/4}.             \tag{CD1}
\]
Moreover
\[
 \mathfrak N_{\rm cen}
 =\sup_{P_jY_j=Y_jP_j=0}
 \sum_j\E\left[I_j\left(
 \frac{2\langle Y_j,X_j\rangle}{m_j^{3/2}}
                     -a_j\|Y_j\|_{\op}^2\right)\right]. \tag{CD2}
\]
The supremum permits every conditionally measurable real matrix
family, with no restriction on efficient computation.
\end{proposition}
\begin{proof}
Suppress $j$. Orthogonal compression is contractive for nuclear norm.
Also
\[
 E'-X=PE'+(I-P)E'P
\]
has rank at most two and Frobenius norm at most $\|E'\|_{\rm F}$:
the three blocks $X$, $PE'$, and $(I-P)E'P$ are Frobenius-orthogonal.
Thus $\|E'-X\|_*^2\le2\|E'\|_{\rm F}^2\le18m^2$.
The nuclear triangle inequality and $(u+v)^2\le2u^2+2v^2$
give
$\|E'\|_*^2\le2\|X\|_*^2+36m^2$.
On the indicated event, $1/(ma)\le2^{k_0-k}\le2^{-k_0/4}$.
There are at most $s$ actual cuts. This proves (CD1).

Write an SVD of $X$ using only its nonzero singular values:
$X=U\Sigma V^{\mathsf T}$. Its left and right singular vectors
are orthogonal to $\boldsymbol1$. Therefore $Z=UV^{\mathsf T}$
is centered, has operator norm one when $X\ne0$, and satisfies
$\langle Z,X\rangle=\|X\|_*$.
Conversely $|\langle Y,X\rangle|\le\|Y\|_{\op}\|X\|_*$.
Maximizing
$2t\|X\|_*/m^{3/2}-at^2$ over $t\ge0$ gives
$\|X\|_*^2/(am^3)$, attained at
$t=\|X\|_*/(am^{3/2})$. Use $Y=0$ when $X=0$.
For fixed $N$ the history space is finite, so the optimizing choices
can be made simultaneously. This justifies the supremum of the sum
and expectation, rather than only a pointwise upper bound.
\end{proof}

The missing assertion is
\[
 \mathfrak N_{\rm cen}\le Cb_Q+
                K M_n^\alpha(s+1)^\beta e^{-cN}.        \tag{CD3}
\]
If proved with fixed explicit constants, (CD1) and Z2 would give
T5 with $C_T=6C$, $K_T=6K+65644$,
$\beta_T=\max(\beta,1)$, $\alpha_T=\alpha$ and
$c_T=\min(c,2^{-137})$. Here $\log2\ge1/2$ bounds the new
remainder. The reviewed choice
$1+1/(768(C_T+1))$ would then apply. No such $C$ is established.

\section{Two copies: the normalization matters}
For a matrix $A$ and two row and column strings define
\[
 \delta A=A(u,v)-A(u',v)-A(u,v')+A(u',v').
\]
If $u,u',v,v'$ are independent uniform strings and $Y$ is centered,
\[
 \E[\delta Y\,\delta E']
           =\frac4{m^2}\langle Y,E'\rangle
           =\frac4{m^2}\langle Y,X\rangle.              \tag{CD4}
\]
This follows by expanding the sixteen terms and using zero row and
column sums. It retains signs. It does not bound the pairing by defect.

Independent complete bit strings have an unsuitable raw union
normalization: each differing bit changes two distinct source groups,
so the expected value of $F$ at their union is $(5/2)^{2q}$.
The following coupling removes this particular exponential factor.
It is not a replacement of the original marginal distribution.

Let $\eta$ have the unchanged original pre-bit law, and let
$b\in\{0,1\}^{2q}$ be its original uniform bit string. Independently
flip each bit with probability
\[
 p=\frac1{2q},\qquad b'=b\mathbin{\oplus}e,\qquad
 U^*=X_{\eta,b}\vee X_{\eta,b'} .
\]
Write $\mu$ for this joint law. For cut $j$, let $e_{j,r}$ and
$e_{j,c}$ retain respectively only its honored row and column flips.
Its four mixed corners have strings
$b$, $b\oplus e_{j,r}$, $b\oplus e_{j,c}$,
$b\oplus e_{j,r}\oplus e_{j,c}$.
Their union $U_j$ is below $U^*$.

\begin{proposition}
All four corners have the original marginal law, also conditionally
on $\eta$ and that cut's unhonored bits. Furthermore
\[
 F(U^*)=4^{|e|},\qquad
 Z_q:=\E_\mu F(U^*)=(1+3/(2q))^{2q}\le e^3.             \tag{CD5}
\]
The probability law $d\widehat\mu=F(U^*)\,d\mu/Z_q$
has the same pre-bit law and the same uniform $b$, independently
of a flip mask with flip probability
\[
                       \widehat p=\frac{4p}{1+3p}.       \tag{CD6}
\]
Thus it too preserves all four original corner marginals.
The actual cuts and the actual A4 mixed contexts obey the one
joint budget
\[
 \begin{split}
 \sum_j\E_\mu D_{Q_j}(U_j)
 &\le \E_\mu D_Q(U^*)\\
 &=\sum_{v\ {\rm product}}\E_\mu H_v(U^*)J_v(U^*)\\
 &=Z_q\,\E_{\widehat\mu}\frac{D_Q(U^*)}{F(U^*)}
 \le Z_q b_Q\le e^3b_Q .                               \tag{CD7}
 \end{split}
\]
\end{proposition}
\begin{proof}
Conditional on $\eta$, the honored sets are fixed. XOR with a mask
independent of uniform bits preserves their joint uniform law.
The mixed corners never change the unhonored bits, so the same
argument holds with those bits fixed.

In the original sampler every bit changes the outgoing choice at
exactly two source groups. Different bits use disjoint groups.
A flipped bit therefore contributes four to $F(U^*)$ and an unflipped
bit contributes one. This proves the first equality of (CD5).
Independence gives its second equality, and $1+x\le e^x$ gives
the bound. The weight $4^{|e|}$ depends on neither $b$ nor $\eta$.
Factoring its normalization bit by bit gives (CD6).
It also proves that the displayed density defines a probability law.

Each $D_{Q_j}$ has nonnegative coefficients after Booleanization.
Hence $D_{Q_j}(U_j)\le D_{Q_j}(U^*)$.
Sum actual cuts at this one common input and use A7.
Use A4 at the same input for the next equality.
Finally $D_Q/F\le b_Q$ everywhere used here.
All unions have every group nonempty, so every division is legitimate.
No separately optimized gate maximum, degree factor, or factor $s$
has been placed on $b_Q$.
\end{proof}

\section{A sparse-noise inverse without a dimension loss}
Centering alone would leave the smallest nonzero noise eigenvalue
$1/q$. Low Walsh degrees can instead be removed at exponentially
small nuclear-energy cost.
Fix a conditional matrix, set $k=k_j$, and assume $I_j=1$.
For admissible $N\ge2^{142}$, $k\ge320$. Set
\[
 r=\lfloor k/256\rfloor,\quad
 d_k=\sum_{\ell=0}^{r-1}\binom{k}{\ell},\quad
 L=\hbox{orthogonal projection onto Walsh degrees less than }r,
 \quad H=I-L.                                          \tag{CD8}
\]
Here $H$ is a matrix on row strings, and the same definition is
used on column strings. These are analysis operators, not circuit gates.
The normalized Walsh vector of a subset $A$ is
$2^{-k/2}(-1)^{\sum_{i\in A}u_i}$.

Let $T$ be the bit-flip transition matrix on these $k$ bits,
with the same $p=1/(2q)$ as above. Its Walsh eigenvalue in degree
$\ell$ is $\rho^\ell$, where $\rho=1-1/q$.
Define $B$ to be zero on degrees below $r$ and to have eigenvalue
$(1-\rho^\ell)^{-1}$ in degree $\ell\ge r$.

\begin{proposition}
Uniformly in the allowed parameters,
\[
 d_k\le2^{k/16},\qquad
 B(I-T)=(I-T)B=H,\qquad \|B\|_{\op}\le2^{121}.           \tag{CD9}
\]
For $X^{\rm hi}_j=H_jE'_jH_j$ and
$\mathfrak N_{\rm hi}=\sum_j\E[I_j\|X^{\rm hi}_j\|_*^2/(m_j^3a_j)]$,
\[
 \mathfrak N_{\rm hi}\le\mathfrak N_{\rm cen}\le\mathfrak N
 \le2\mathfrak N_{\rm hi}+36s\,2^{-11k_0/64}.            \tag{CD10}
\]
\end{proposition}
\begin{proof}
If $r=1$, then $d_k=1$. Otherwise put $h=r-1\le k/256$.
At $z=h/k\le1$,
$z^h\sum_{\ell=0}^h\binom{k}{\ell}\le(1+z)^k\le e^h$.
The increasing function $x\log(ek/x)$ on $(0,k]$ gives
\[
 d_k\le(ek/h)^h\le(256e)^{k/256}
       \le2^{10k/256}\le2^{k/16}.
\]
The Walsh eigenvalue formula follows by independence of the flips.
Since $r\ge k/512$, $r\le q$, and
$(1-1/q)^r\le e^{-r/q}\le(1+r/q)^{-1}$, the eigenvalues
of $I-T$ on $H$ are at least
\[
 \frac{r}{q+r}\ge\frac{k}{1024q}
 \ge\frac{k_0}{1024q}\ge2^{-121}.
\]
The last step uses $q\le N/2^{23}=2^{111}k_0$.
This proves (CD9).

The matrix $E'-HE'H=LE'+HE'L$ has rank at most $2d_k$
and Frobenius norm at most $\|E'\|_{\rm F}$, by orthogonal blocks.
Exactly as in (CD1),
$\|E'\|_*^2\le2\|HE'H\|_*^2+36d_km^2$.
On the indicated event,
\[
 \frac{d_k}{ma}\le2^{k_0-15k/16}
                         \le2^{-11k_0/64}.
\]
Compression gives the other inequalities, since $H(I-P)=H$.
Sum only over the at most $s$ original cuts.
\end{proof}

For arbitrary real $Y$, let $\widetilde Y=BYB$. Under the
conditional sparse-noise law, expansion in Walsh vectors, or
twice the identity
$\E[(f(u)-f(u'))(g(u)-g(u'))]
 =2\langle f,(I-T)g\rangle/m$, gives
\[
 \begin{split}
 \E_e\E_{u,v}[\delta\widetilde Y\,\delta E']
  &=\frac4{m^2}\langle\widetilde Y,(I-T)E'(I-T)\rangle\\
  &=\frac4{m^2}\langle Y,HE'H\rangle,\qquad
 \|\widetilde Y\|_{\op}\le2^{242}\|Y\|_{\op}.             \tag{CD11}
 \end{split}
\]
All expectations in this identity use $\mu$, not the tilted law.
The tilt was used only to justify (CD7).
The high-frequency energy is therefore exactly the supremum of
\[
 \sum_j\E_{\eta,z_j}\left[I_j\left(
  \frac{\sqrt{m_j}}2
       \E_{u,v,e}[\delta(B_jY_jB_j)\,\delta E'_j]
              -a_j\|Y_j\|_{\op}^2\right)\right],         \tag{CD12}
\]
where one may restrict to $Y_j=H_jY_jH_j$.
The SVD argument for (CD2) applies verbatim.
A fixed defect bound for (CD12) would also yield T5 with coefficient
$6C$; its new remainder rate is at least $11/2^{141}$.
Neither (CD7) nor (CD11) supplies that bound.

\section{Why the common budget is not yet a signed charge}
The entries in $\delta E'_j$ use the fixed \emph{global} completion
at $U_\eta$. In (CD7), the local union $U_j$ usually contains
only a few additional choices. Replacing the former completion by
the latter changes the matrix. The following exact diagnostic
shows that even a pointwise zero-defect implication fails.
It does not refute an averaged estimate with an exponential remainder.

Use three abstract gadgets, with twelve two-variable groups
$(z_i,w_i)$, $0\le i<12$. Row bit $u_t$ selects the same
option in groups $2t,2t+1$; column bit $v_t$ does so in
groups $6+2t,7+2t$. Let $S=\{0,\ldots,5\}$, and put
\[
 \begin{split}
 f_S&=(z_0+2w_0^3)\prod_{i=1}^{5}(z_i+w_i),\\
 V&=(z_8+2w_8^2)
                  \prod_{i\in\{6,7,9,10,11\}}(z_i+w_i),\\
 A&=\tfrac12 f_S(1+\tfrac12w_6),\qquad Q=AV.            \tag{CD13}
 \end{split}
\]
Build $V$ before $A$ and their final product. The last balanced
gate is $A$, $G(A)=S$, and the divided context is exactly $V$.
Cutting $A$ annihilates the output. Thus this is one actual cut,
not a freely chosen product representation. Repeated powers and
shared variable inputs are retained.

On the full Boolean cube, $Q\le3F$. After Booleanization,
\[
 D_Q=\tfrac14\,\overline f_S\,w_6z_6\,\overline V_{\ne6},
 \qquad b_Q=\tfrac12.                                  \tag{CD14}
\]
Bars mean idempotent powers; $V_{\ne6}$ omits group six.
Indeed $\overline f_S/F_S\le2$ and
$\overline V_{\ne6}/F_{B\setminus\{6\}}\le2$.
If $w_6z_6=1$ its normalized factor is $1/2$.
All three maxima occur together, proving the exact defect value.
Full-group coverage handles the zero faces.

In the complete original-bit box of this abstract restriction,
with $f(u)=1+u_0$ and $h(v)=(1+v_1)(1-v_0)$,
\[
 \begin{split}
 M(u,v)&=\tfrac12(1+u_0)(1+\tfrac12v_0)(1+v_1),\\
 R(u,v)&=\tfrac34(1+u_0)(1+v_1)\le3,\\
 E'(u,v)&=\tfrac14 f(u)h(v),\\
 X(u,v)&=\tfrac14(f(u)-\tfrac32)(h(v)-\tfrac34).
                                                               \tag{CD15}
 \end{split}
\]
At the finite heavy cutoff zero, the heavy mean is $a=9/32$.
For the finite parameter $k_0=2$, $k=3\ge5k_0/4$ and
$a\ge2^{-k_0}$. These finite parameters are not the admissible
asymptotic recipe for $N$.

The eight-square Gram matrix $XX^{\mathsf T}$ has exact spectrum
\[
           (11/16)^1,\quad 0^7,\qquad
 \frac{\|X\|_*^2}{m^3a}=\frac{11}{2304}.                \tag{CD16}
\]
For example $\sum_u(f(u)-3/2)^2=2$ and
$\sum_v(h(v)-3/4)^2=11/2$, so the squared rank-one norm is
$11/16$. This is a genuine centered error, not a row-only proxy.

Take $u=0$, $u'=1$, $v=0$, $v'=2$, with bit indices starting
at zero. Thus only row bit zero and column bit one flip.
The optional variable $w_6$ is absent at all four corners and
at their union $U^*$. Nevertheless, for the centered choice $Y=X$,
\[
 D_Q(U^*)=0,\quad F(U^*)=16,\quad
 \delta E'=\delta X=\tfrac14,\quad
                          \delta X\,\delta E'=\tfrac1{16}. \tag{CD17}
\]
For $p=1/6$ and the finite cutoff $r=1$, the same example gives
$\delta(BXB)=9/5$ and
$\delta(BXB)\delta E'=9/20>0$.
Consequently zero local defect does not force the signed integrand
in (CD12) to vanish. Applying A4 to that integrand without another
identity or inequality is not justified.
The example has rank one and can be absorbed by an exponential
low-rank remainder as dimensions grow. It supplies no all-size
counterexample to (CD3), Z3, T5, or the target lower bound.

\section{The pointwise shortcut also fails with the actual heavy threshold}
The preceding finite example alone does not address the asymptotic
cutoff $h=\lfloor N/2^{150}\rfloor$. The following construction does.
It rules out only a remainder-free pointwise use of (CD7), not the
averaged charge (CD12).

Fix any pre-bit history in the support of the original sampler.
Choose a balanced $S$ containing all its $2q$ row sources and none
of its $2q$ column sources. There are enough other groups to fill
out $S$, since $q\le N/2^{23}$. The circuit below is then fixed;
the original sampler is not conditioned or changed in the target
claim. At this particular history, all $k=q$ gadgets are honored.
Set
\[
 d=h+1,\quad D=\{0,\ldots,d-1\},\quad
 W=\{d,\ldots,k-1\}.
\]
One has $d+1\le k_0$, $d<r=\lfloor k/256\rfloor$, and
$k-d\ge r$. Indeed $k_0\ge256$,
$d\le k_0/2^{16}+1$, and
$r\ge k/512\ge N/2^{33}>k_0$.

On one representative source of each row bit, let
$P_{\rm even,S}$ select the even-parity choices of its two
designated arcs, and multiply by the full group sums on all
other groups of $S$. Define $P_{\rm even,B}$ similarly on one
representative source for each column bit in $W$, with full group
sums on the remaining groups of $B=\overline S$.
Both polynomials are coefficientwise at most the respective full
group product. They have positive linear-size circuits: starting
from $E=1,O=0$, a designated pair $(z,w)$ updates the two states to
$(Ez+Ow,Oz+Ew)$. The final even state is multiplied by the
remaining group sums.

For $i\in D$, let $y_i$ be the bit-one nonroot arc at its
representative column source, and put $C_D=\prod_{i\in D}y_i$.
Define
\[
 f_S=F_S+P_{\rm even,S},\quad
 V=F_B+P_{\rm even,B},\quad
 A=\tfrac12 f_S(1+\tfrac12 C_D),\qquad Q=AV.             \tag{CD18}
\]
Again build $V$ before $A$. This is an $O(M_n+q+d)$-gate
original DAG with exactly one nonzero actual cut, whose context
is $V$. It has full-group coverage and $Q\le3F$ on the entire
Boolean cube. Precisely the branch $\tfrac14f_SC_DV$ is heavy:
its $d=h+1$ optional nonroot source groups lie outside $S$;
the other branch and $V$ have no optional group in the opposite side.

Write $\chi_U$ for parity on all row bits, $\chi_W$ for parity
on the indicated column bits, and
$c(v)=1_{\{v_i=1\ (i\in D)\}}$. At the chosen history,
\[
 f(u)=\tfrac32+\tfrac12\chi_U(u),\quad
 g(v)=\tfrac32+\tfrac12\chi_W(v),\quad
 a=\tfrac9{16}2^{-d}>2^{-(d+1)}\ge2^{-k_0}.
\]
The original completion is $R=\tfrac34fg\le3$, so
$E'=\tfrac14fg(1-c)$ without clipping.
The low column degrees of $1-c$ are at most $d<r$; multiplying
by $\chi_W$ gives only degrees $k-d,\ldots,k$.
Consequently
\[
 X^{\rm hi}(u,v)=\frac{\chi_U(u)\chi_W(v)(1-c(v))}{16},
 \qquad
 \|X^{\rm hi}\|_{\op}^2
          =\frac{m^2(1-2^{-d})}{256}.                    \tag{CD19}
\]
This follows directly from the two rank-one factor norms.

Take the all-zero row and column strings, flip one row bit and
one column bit in $W$, and flip no others. All $y_i$ are absent
from their common union, so $D_Q(U^*)=0$. Nevertheless
$\delta E'=1/4$. Set $\lambda_\ell=1-\rho^\ell$. At these strings
the column inverse of $\chi_W(1-c)$ has values $J$ and $-J$, where
\[
 \begin{split}
 J&=\frac1{\lambda_{k-d}}
   -2^{-d}\sum_{t=0}^d
        \frac{(-1)^t\binom dt}{\lambda_{k-d+t}}\\
  &=\sum_{n=0}^{\infty}\rho^{n(k-d)}
                    \bigl[1-2^{-d}(1-\rho^n)^d\bigr]\ge1 .
 \end{split}
\]
The geometric series converge since $k-d>0$ and $0\le\rho<1$.
Every summand is nonnegative, and the term $n=0$ is one.
Thus
$\delta(BX^{\rm hi}B)\delta E'=J/(16\lambda_k)>0$.
Use the admissible, conditionally fixed high-frequency test matrix
$Y=m^{-2}X^{\rm hi}$. Since $a\le9/32$,
\[
 \begin{split}
 \frac{\sqrt m}{2}\delta(BYB)\delta E'-a\|Y\|_{\op}^2
 &=\frac{J}{32\lambda_km^{3/2}}
             -\frac{a(1-2^{-d})}{256m^2}\\
 &\ge\frac{247}{8192m^2}>0,\qquad D_Q(U^*)=0.            \tag{CD20}
 \end{split}
\]
This event has positive probability in the full original law
and in the sparse flip law; no lower bound on that probability
is asserted.
For every finite $C$, it refutes the proposed pointwise estimate
that the sum of these regularized signed integrands is at most
$C D_Q(U^*)$ with no remainder.

It does not refute the actual mission. At \emph{every} history,
$f_S$ depends only on the row side, while $V$ and $C_D$ depend
only on the column side after the unhonored bits are fixed.
The unclipped completion is always at most three, and
$0\le E'\le1$ has rank at most one. Hence this entire original-sampler
family satisfies
$\mathfrak N\le2^{-k_0/4}$ on the strong-good/thick event.
An allowed exponential remainder absorbs it.
The remainder-free pointwise route is retired for (CD20);
the averaged signed route remains open.

\paragraph{Exact replay and its scope.}
\path{centered_dual_check.py} binds (CD15)--(CD17) to the original
rational DAG, its actual cut, exact coefficient-preserving projection,
and heavy split. It checks full-cube domination through all 64 inputs
on the three distinguished groups and exact cancellation of the
other group factors, rather than claiming individual enumeration of
the entire physical cube.
It checks all 4096 three-gadget base-string/flip-mask outcomes,
the original and tilted corner marginals for every honored subset,
the common union normalization and defect budget, and sparse-noise
polarization with two finite Walsh cutoffs. Two further finite
Walsh transforms check the high-frequency formulas and inverse
in (CD19)--(CD20); they do not simulate an admissible large graph.
The general proofs above, not these finite cases, supply the
all-size statements.

The new bounded Jacobian operation
\texttt{matrix.symmetric.rational\_spectrum\_claim.check} checks
the complete spectrum of the circuit-bound centered Gram matrix.
Its inspected contract, exact payload and full typed result are in
the dedicated \path{centered-spectrum-*.json} files. The replay
independently checks every shifted nullity and binds the returned
matrix and claims to the circuit. No numerical eigenvalue or timeout
is used as evidence.

\paragraph{Remaining obligation.}
The construction now has original marginals, a bounded normalization,
one whole-output A4/A7 budget, and a bounded high-frequency inverse.
What remains is an inequality controlling the supremum (CD12) by
that budget with its actual heavy denominator and one exponentially
small original-size remainder. The global-to-local completion
comparison and the signed factor $\sqrt m\,\delta(BYB)$ are
uncontrolled. No summation of absolute component budgets, efficient
semantic extraction, nonnegative-coordinate Gram restriction, or
entrywise monotonicity of nuclear norm is substituted for this step.
It remains a mathematical gap, not an external-tool blocker.
\end{document}
